\scanpage{56}
\mspage{4 - 36 -}
n\textsuperscript{o} 392

Revenant à la démonstration de la prop. 19, on a, en vertu du lemme 4, $\handwritten{\bigcap}_{\underline P\in P(B)}R_{\underline P}=\handwritten{\bigcap}_{\underline p\in P(A)}(R_{[A]})_{\underline p}$, et la conclusion résulte du n\textsuperscript{o} 2, th. 2.

PROPOSITION 20.- (i) \uline{Pour qu'un} $A$-\uline{module de type fini} $M$ \uline{soit pseudo-nul}, \uline{il faut et il suffit que} $M\otimes_A B$ \uline{soit un} $B$-\uline{module pseudo-nul}.

(ii) \uline{Si} $M$ \uline{est un} $A$-\uline{module de torsion de type fini}, $M\otimes_A B$ \uline{est un} $B$-\uline{module de torsion de type fini}, \uline{et l'on a} :
\[
\handwritten{\chi}_B(M\otimes_A B)=i(\handwritten{\chi}_A(M)).\tag*{\makebox[0pt][r]{\makebox[\displaywidth][l]{(22)}}}
\]

(iii) \uline{Si} $M$ \uline{est un} $A$-\uline{module de type fini}, $M\otimes_A B$ \uline{est un} $B$-\uline{module de type fini}, \uline{et l'on a} :
\[
c_B(M\otimes_A B)=i(c_A(M))\tag*{\makebox[0pt][r]{\makebox[\displaywidth][l]{(23)}}}
\]
\[
r_B(M\otimes_A B)=r_A(M).\tag*{\makebox[0pt][r]{\makebox[\displaywidth][l]{(24)}}}
\]

(i) Soient $\underline P$ un idéal premier de $B$, $\underline p=\underline P\handwritten{\cap} A$ ; on a $(M\otimes_A B)_{\underline P}=M\otimes_A B_{\underline P}$, (chap. II, \S~2, n\textsuperscript{o} 7, prop. 18), et d'autre part $M\otimes_A B_{\underline P}=(M\otimes_A A_{\underline p})\otimes_{A_{\underline p}}B_{\underline P}={}=M_{\underline p}\otimes_{A_{\underline p}}B_{\underline P}$ ; la relation $M_{\underline p}=0$ est donc équivalente à $(M\otimes_A B)_{\underline P}=0$, (chap. II, \S~4, n\textsuperscript{o} 4, lemme 4). Il suffit d'appliquer cette remarque à l'idéal $\underline P=(0)$ et aux idéaux $\underline P\in P(B)$ pour prouver (i), compte tenu du n\textsuperscript{o} 4, prop. 9.

Pour prouver (ii), nous utiliserons le lemme suivant :

LEMME 5.- \uline{Soient} $M_1$, $M_2$ \uline{deux} $A$-\uline{modules de type fini}, $f:M_1\longrightarrow M_2$ \uline{un homomorphisme injectif}. \uline{Alors le noyau de} $f\otimes 1_B:M_1\otimes_A B\longrightarrow M_2\otimes_A B$ \uline{est pseudo-nul}.

Soit $\underline p$ un idéal premier de $A$ de hauteur $\leq 1$. On a $(M_i\otimes_A B)_{\underline p}=(M_i)_{\underline p}\otimes_{A_{\underline p}}B_{\underline p}$ $(i=1,2)$, (chap. II, \S~2, n\textsuperscript{o} 7, prop. 18), et $(f\otimes 1_B)_{\underline p}=f_{\underline p}\otimes 1_{B_{\underline p}}$ ; l'hypothèse que $f$ est injectif entraîne qu'il en est de même de $f_{\underline p}$, (chap. II, \S~1, n\textsuperscript{o} 4, th. 1) ; d'autre part, vu le choix de $\underline p$, $A_{\underline p}$ est un anneau principal et $B_{\underline p}$ un $A_{\underline p}$-module sans torsion de type fini, donc \uline{libre} ; on en conclut que $f_{\underline p}\otimes 1_{B_{\underline p}}$ est
% archive sheet number: 55
